\documentclass{article}
\usepackage{graphicx} % Required for inserting images
\usepackage{amsmath, amssymb, amsthm} % Math stuff
\usepackage[letterpaper, top=1.0 in, bottom=1.0in, left = 1.0in, right=1.0in, heightrounded]{geometry}
\usepackage{setspace}
\theoremstyle{definition}
\onehalfspacing

\title{An Elegant Solution}
\author{Noah Park}
\date{April 2026}

\begin{document}

\maketitle

\section{The Problem}

Find the positive root of $\frac {1}{x^2-10x-29} + \frac{1}{x^2-10x-45} - \frac {2}{x^2-10x-69} = 0$.

\section{The Solution}
\begin{align*}
\frac {1}{x^2-10x-29} + \frac{1}{x^2-10x-45} - \frac {2}{x^2-10x-69} &= 0\\
\frac {1}{x^2-10x-29} + \frac{1}{x^2-10x-45} &=  \frac {2}{x^2-10x-69}
\end{align*}


Let $y = x^2-10x-37$.

\[
\frac {1}{y+8} + \frac{1}{y-8} =  \frac {2}{y-32}
\]

Combine the left-hand side:
\begin{align*}
\frac {y-8}{y^2-64} + \frac{y+8}{y^2-64} &=  \frac {2}{y-32}\\
\frac {2y}{y^2-64} &=  \frac {2}{y-32} \\
\frac {y}{y^2-64} &=  \frac {1}{y-32}
\end{align*}

Cross-multiplying:
\begin{align*}
y^2-32y&=y^2-64 \\
-32y&=-64 \\
y&=2
\end{align*}

Substituting back:
\begin{align*}
x^2-10x-37&=2\\
x^2-10x-39&=0\\
(x-13)(x+3)&=0
\end{align*}

Thus,
\[
x= \boxed{13}
\]

\end{document}
